\section{Discusión}

Los resultados muestran que el tipo de cambio UYU/BRL presenta variación apreciable durante el período estudiado. La media muestral de \MediaCambio{} y el desvío estándar de \DesvioCambio{} ofrecen una primera síntesis del comportamiento central y de la dispersión, pero la serie temporal evidencia cambios de nivel y períodos de distinta volatilidad.

El análisis gráfico indica que la distribución empírica de la cotización no debe interpretarse aisladamente de su evolución temporal. Una cotización financiera observada durante muchos años puede mezclar distintos contextos económicos, por lo que una única distribución teórica puede ser útil como aproximación descriptiva, pero no necesariamente como modelo completo del fenómeno.

Los log-retornos ayudan a estudiar las variaciones diarias de forma más estable. Aun así, la normalidad no puede asumirse sin revisión, y la autocorrelación debe considerarse al interpretar intervalos y pruebas. En este sentido, los métodos clásicos son didácticamente útiles y permiten aplicar herramientas del curso, pero deben complementarse con enfoques empíricos o específicos de series temporales cuando se busca mayor rigurosidad.

Una limitación importante es que los datos no provienen de un muestreo aleatorio simple. Se trata de observaciones históricas disponibles, condicionadas por días de mercado, disponibilidad de cotizaciones y el procedimiento de construcción indirecta de UYU/BRL. Además, las asociaciones entre variables cambiarias no implican causalidad: pueden existir factores externos comunes que afecten simultáneamente al peso uruguayo, al real brasileño y al dólar estadounidense.

